% from ega4-15-terminal-fr.tex
\oldpage[IV-3]{254}
\EGAIVthreeContentsLine{9.4.}{Constructibilité de certaines propriétés des Modules}{62}
\EGAIVthreeContentsLine{9.5.}{Constructibilité de propriétés topologiques}{67}
\EGAIVthreeContentsLine{9.6.}{Constructibilité de certaines propriétés des morphismes}{71}
\EGAIVthreeContentsLine{9.7.}{Constructibilité des propriétés de séparabilité, d'irréductibilité géométrique et de connexité géométrique}{76}
\EGAIVthreeContentsLine{9.8.}{Décomposition primaire au voisinage d'une fibre générique}{83}
\EGAIVthreeContentsLine{9.9.}{Constructibilité des propriétés locales des fibres}{88}

\EGAIVthreeContentsLine{\S~10.}{Préschémas de Jacobson}{95}
\EGAIVthreeContentsLine{10.1.}{Parties très denses d'un espace topologique}{95}
\EGAIVthreeContentsLine{10.2.}{Quasi-homéomorphismes}{97}
\EGAIVthreeContentsLine{10.3.}{Espaces de Jacobson}{101}
\EGAIVthreeContentsLine{10.4.}{Préschémas de Jacobson et anneaux de Jacobson}{101}
\EGAIVthreeContentsLine{10.5.}{Préschémas de Jacobson noethériens}{104}
\EGAIVthreeContentsLine{10.6.}{Dimension dans les préschémas de Jacobson}{107}
\EGAIVthreeContentsLine{10.7.}{Exemples et contre-exemples}{109}
\EGAIVthreeContentsLine{10.8.}{Profondeur rectifiée}{110}
\EGAIVthreeContentsLine{10.9.}{Spectres maximaux et ultrapréschémas}{112}
\EGAIVthreeContentsLine{10.10.}{Espaces algébriques de Serre}{114}

\EGAIVthreeContentsLine{\S~11.}{Propriétés topologiques des morphismes plats de présentation finie. Critères de platitude}{116}
\EGAIVthreeContentsLine{11.1.}{Ensembles de platitude (cas noethérien)}{117}
\EGAIVthreeContentsLine{11.2.}{Platitude d'une limite projective de préschémas}{119}
\EGAIVthreeContentsLine{11.3.}{Application à l'élimination d'hypothèses noethériennes}{132}
\EGAIVthreeContentsLine{11.4.}{Descente de la platitude par des morphismes quelconques : cas d'un préschéma de base artinien}{143}
\EGAIVthreeContentsLine{11.5.}{Descente de la platitude par des morphismes quelconques : cas général}{150}
\EGAIVthreeContentsLine{11.6.}{Descente de la platitude par des morphismes quelconques : cas d'un préschéma de base unibranche}{154}
\EGAIVthreeContentsLine{11.7.}{Contre-exemples}{157}
\EGAIVthreeContentsLine{11.8.}{Un critère valuatif de platitude}{159}
\EGAIVthreeContentsLine{11.9.}{Familles séparantes et universellement séparantes d'homomorphismes de faisceaux de modules}{160}
\EGAIVthreeContentsLine{11.10.}{Familles schématiquement dominantes de morphismes et familles schématiquement denses de sous-préschémas}{170}

\EGAIVthreeContentsLine{\S~12.}{Étude des fibres des morphismes plats de présentation finie}{173}
\EGAIVthreeContentsLine{12.0.}{Introduction}{173}
\EGAIVthreeContentsLine{12.1.}{Propriétés locales des fibres d'un morphisme plat localement de présentation finie}{174}

\clearpage
